% Nesse arquivos estão as importações de todos os pacotes usados (e que possívelmente serão usados por um aluno).

% Recomendo que você remova os que não vai usar pelo bem do tempo de compilação.

%========================
% Idioma
%========================
\usepackage[brazil]{babel}

%========================
% Fontes e tipografia
%========================
\usepackage{fontspec}
\usepackage{textcomp}

%========================
% Layout e formatação
%========================
\usepackage{geometry}         % Margens da página
\usepackage{setspace}         % Espaçamento entre linhas
\usepackage{parskip}          % Espaço entre parágrafos
\usepackage{indentfirst}      % Indenta o primeiro parágrafo
\usepackage{setspace}
\usepackage{titlesec}         % Personalização de títulos

%========================
% Gráficos e figuras
%========================
\usepackage{graphicx}         % Inserção de imagens
\usepackage{caption}          % Customização de legendas
\usepackage{subcaption}       % Subfiguras
\usepackage{float}            % Controle de posição de figuras
\usepackage{tikz}             % Desenhos e diagramas


%========================
% Matemática
%========================
\usepackage{amsmath}          % Ambiente matemático avançado
\usepackage{amssymb}          % Símbolos matemáticos
\usepackage{cancel}           % Cancelamento em equações
\usepackage{commath}          % Comandos matemáticos úteis
\usepackage{physics}          % Notação de física

%========================
% Química
%========================
\usepackage{chemformula}      % Fórmulas químicas

%========================
% Listas e tabelas
%========================
\usepackage{enumitem}         % Customização de listas
\usepackage{booktabs}         % Tabelas profissionais
\usepackage{multicol}         % Múltiplas colunas
\usepackage{multirow}

%========================
% Referências e links
%========================
\usepackage{hyperref}         % Hiperlinks
\usepackage{url}              % URLs
\usepackage[alf]{abntex2cite}
% \citebrackets[]{}


%========================
% Código-fonte
%========================
\usepackage[newfloat]{minted} % Inserção de código (requer -shell-escape)

%========================
% Elementos adicionais
%========================
\usepackage{pdfpages}         % Inserir PDFs
\usepackage[absolute,overlay]{textpos} % Posicionamento absoluto de texto
\usepackage{color,soul}
\usepackage{lipsum}
\usepackage{etoolbox}